%BAB III: HARDSKILL - PENGEMBANGAN FITUR DAN MODUL PROYEK
\section{MBKM Hardskill: Pengembangan Fitur dan Modul Proyek}

\begin{table}[H]
    \centering
    \begin{tabular}{ll}
        \textbf{Nama} & : Muhammad Argya Vityasy \\
        \textbf{NIM} & : 23/522547/PA/22475 \\
        \textbf{Pembimbing Akademik} & : Guntur Budi Herwanto, S.Kom., M.Cs. \\
        \textbf{Pembimbing Program MBKM} & : Guntur Budi Herwanto, S.Kom., M.Cs. \\
        \textbf{Mata Kuliah} & : Internship: Pengembangan Fitur dan Modul Proyek \\
        \textbf{Kode MK} & : MII21-3011 \\
        \textbf{Total SKS} & : 4 SKS \\
        \textbf{Judul Magang} & : MBKM Mandiri GDP Labs: Digital Employee \\
        \textbf{Durasi} & : 9 Februari 2026 -- 22 Juni 2026 \\
    \end{tabular}
\end{table}

\subsection{Landasan Teori}

Digital Employee (DE) adalah platform agen otomasi berbasis AI yang dikembangkan GDP Labs. Agen berkomunikasi dengan layanan eksternal melalui Model Context Protocol (MCP) dan setiap tool wajib memenuhi kontrak \texttt{args\_schema} (Pydantic BaseModel), \texttt{tool\_config\_schema}, \texttt{get\_tool\_config()}, serta metode \texttt{\_run()}. Kontrak ini membuat tool modular dan dapat diuji secara independen. Pola coordinator-subagent dipakai untuk mendelegasikan tugas dari agen utama ke sub-agen yang lebih spesifik berdasarkan intensi pengguna.

Kegiatan pengembangan saya pada mata kuliah ini berfokus pada tiga aplikasi dalam monorepo \texttt{digital-employee}: \texttt{digital-employee-weekly-report} (DE-WR), \texttt{digital-employee-weekly-report-issue-reoccuring} (DE-WRI), dan \texttt{digital-employee-pm} (DE-PM). Seluruh kode dikirim melalui pull request (PR) di GitHub dan di-review oleh rekan tim sebelum di-merge.

\subsection{Kegiatan}

\subsubsection{DE-WRI: Pipeline Eskalasi Isu Berulang dari Nol}

Kontribusi terbesar pada mata kuliah ini adalah membangun pipeline deteksi dan eskalasi isu berulang (recurring issue) untuk laporan mingguan, yang kemudian di-deploy ke runtime AIP (AI Intelligence Platform) sebagai agen terjadwal. Pipeline ini saya bangun dari nol dalam beberapa langkah iteratif.

\begin{itemize}
    \item \textbf{Inisialisasi Proyek}: Pipeline dimulai sebagai hello world berbasis \texttt{gllm\_pipeline.Pipeline}, lalu ditingkatkan menjadi \textit{deployable pipeline agent} yang dibungkus dalam \texttt{BaseTool} LangChain. TypedDict diganti dengan Pydantic BaseModel setelah muncul error runtime pada server AIP (\texttt{State type must be a TypedDict or Pydantic BaseModel}).
    \item \textbf{FetchEmployeeIssuesTool}: Tool ini melakukan query SQL JOIN \texttt{dummy\_weekly\_reports} $\rightarrow$ \texttt{employees} $\rightarrow$ \texttt{organizations} dan mengelompokkan hasil per pegawai beserta isunya. Pada tahap ini saya juga melakukan migrasi dari \texttt{BosaConnector} ke GLConnectors SDK. Keamanan ditangani dengan validasi nama tabel dan pencegahan SQL injection. Tool ini dilengkapi 56 unit test dan 5 integration test dengan cakupan 100\%.
    \item \textbf{RecurringIssueDetector}: Deteksi isu berulang menggunakan \textit{strategy pattern} dengan \texttt{ExactMatcher} yang mencocokkan string secara case-insensitive, sehingga matcher lain (misalnya \texttt{FuzzyMatcher}) dapat ditambahkan tanpa mengubah alur utama. Satu bug penting diperbaiki di sini: \texttt{issue\_id} semula tidak dimasukkan ke tuple pengelompokan mingguan, sehingga isu yang sama bisa terdeteksi salah sebagai isu baru.
    \item \textbf{GitHub Issue Tools}: Tool GitHub yang semula monolitik direfaktor menjadi tiga tool terpisah, yaitu \texttt{SearchGitHubIssueTool}, \texttt{CommentGitHubIssueTool}, dan \texttt{CloseGitHubIssueTool}, dengan \texttt{LookupGitHubDataTool} untuk memetakan pegawai ke akun GitHub. Sebelum mengeskalasi isu yang sudah tercatat di basis data, status isu di GitHub dicek terlebih dahulu.
    \item \textbf{SendEmailTool}: Tool pengiriman email yang terintegrasi dengan MCP \texttt{google\_mail} BOSA, lengkap dengan unit test dan integration test.
    \item \textbf{EscalationOrchestratorTool}: Pipeline lima langkah (ambil isu, deteksi berulang, cek GitHub, kirim email, catat basis data) dikonsolidasikan menjadi satu tool orkestrasi. \texttt{get\_tool\_config} diseragamkan di seluruh tool.
    \item \textbf{Manajemen Penjadwalan}: Saya menambahkan dukungan manajemen jadwal AIP melalui perintah CLI (buat, hapus, dan info jadwal). Manajemen jadwal dibuat idempoten untuk mencegah bug overwrite, dan \textit{trigger phrase} ditegakkan agar eksekusi terjadwal tidak dijalankan secara tidak sengaja. Jadwal mingguan dikonfigurasi setiap Minggu pukul 10.00.
    \item \textbf{Peningkatan Arsitektur}: Diterapkan pola coordinator dengan dua sub-agen, yaitu Escalation Agent dan GitHub DB Sync Agent. \texttt{GLBaseTool} diganti fungsi utilitas, konfigurasi dipusatkan, fungsi basis data disatukan dengan sanitasi identifier, dan query SQL dieksternalisasi ke berkas \texttt{.sql} lalu di-embed ulang agar kompatibel dengan AIP.
\end{itemize}

Hasil akhirnya adalah agen DE-WRI yang berjalan terjadwal setiap minggu. Keputusan untuk setiap isu berulang mengikuti alur berikut:

\begin{lstlisting}[caption={Decision tree untuk isu berulang}]
# Untuk setiap isu berulang:
# - Cari catatan eskalasi di DB
# - Tidak ditemukan -> Buat isu GitHub baru, catat di DB
# - Ditemukan + status=open -> Verifikasi di GitHub API
#   - GitHub says closed -> Tutup DB record, buat isu baru
#   - GitHub says open -> Tambahkan komentar pada isu yang ada
# - Ditemukan + status=closed -> Tutup DB record, buat isu baru
\end{lstlisting}

Pipeline ini dirangkum dalam PR \#573 (\textit{feat(wr-issue-reoccurrence): implement Weekly Report Escalation Agent with GitHub sync and CLI}).

\subsubsection{DE-WR: Tool dan Pipeline Laporan Mingguan}

Pada aplikasi \texttt{digital-employee-weekly-report}, kontribusi saya mencakup tool, sub-agen, dan pipeline komplain laporan kosong:

\begin{itemize}
    \item \textbf{ReportUpdaterDateTool}: Menambahkan tool \texttt{ReportUpdaterDateTool} dan merefaktor agen \texttt{report\_updater} agar perhitungan tanggal dikerjakan di sisi server, bukan di dalam instruksi agen (PR \#593).
    \item \textbf{Report Email Sender Agent}: Menambahkan sub-agen \texttt{report\_email\_sender\_agent} untuk mengirim email terkait laporan mingguan (PR \#600).
    \item \textbf{Issue Processing Agent}: Menyatukan ekstraksi dan scoring isu ke dalam satu agen \texttt{issue\_processing} agar alur lebih sederhana dan konsisten (PR \#730).
    \item \textbf{Empty Weekly Report Compliance Pipeline}: Membangun pipeline yang mendeteksi laporan mingguan kosong dengan menghitung placeholder \texttt{(please insert here)} pada dokumen Google Docs, mengunggah snapshot PDF ke Google Drive, dan melacak laporan kosong di spreadsheet master (PR \#663). Pipeline ini memiliki \textit{batch run tracker} agar proses bisa dilanjutkan setelah terputus, menamai PDF dengan zona waktu Asia/Jakarta, dan melewati pegawai yang tidak aktif atau tanpa manager (PR \#776). Tool \texttt{google\_drive\_search} yang mendukungnya saya lengkapi dengan test suite hingga cakupan 100\%.
    \item \textbf{Weekly Compliance Summary Email}: Menambahkan \texttt{SendComplianceSummaryEmailTool} untuk mengirim ringkasan komplain mingguan secara otomatis, lengkap dengan auto-resolution dan filter penerima (PR \#714).
\end{itemize}

\subsubsection{DE-WR: E2B Issue Processing}

Sebagai alternatif scoring berbasis LLM langsung, saya mengerjakan jalur E2B yang memproses isu di dalam sandbox Python terisolasi:

\begin{lstlisting}[caption={Alur E2B Issue Processing dalam Weekly Report}]
# Jalur paralel dengan ThreadPoolExecutor
# Setiap batch mendapat sandbox E2B terpisah
# Mengunggah runner_source.py (~800 baris) ke sandbox
# Menginstal dependensi: openai, httpx, python-dotenv, requests
# Menjalankan scoring dengan LLM via OpenRouter
# Menangani PII anonymization via NER API
# Timeout per tool call: 7200 detik
\end{lstlisting}

\subsubsection{DE-PM: E2B Analytics Agent untuk Laporan MoM Mingguan}

Pada aplikasi \texttt{digital-employee-pm}, saya membangun \texttt{E2B Analytics Agent} untuk membuat laporan Minutes of Meeting (MoM) mingguan secara otomatis. Data utama diambil dari Meemo, dengan kredensial yang diturunkan dari \texttt{MEEMO\_MCP\_X\_API\_KEY} dan \texttt{MEEMO\_MCP\_URL} (PR \#2167, dilanjutkan pada PR \#830).

\begin{itemize}
    \item \textbf{Generasi Laporan PDF}: Laporan dibuat bergaya konsultan Big 4 menggunakan ReportLab Platypus, dengan grafik Matplotlib untuk tren bulanan dan perbandingan MoM. Email dikirim secara prompt-driven, PDF diberi nama ber-versi dengan tautan Drive, dan dibagikan ke penerima terkait.
    \item \textbf{Timeout dan Resiliensi}: Default timeout AIP 60 detik dilampaui dengan meneruskan \texttt{request\_timeout} ke pemanggilan SDK E2B dan menyediakan override resiliensi \texttt{tool\_timeout\_seconds} melalui \texttt{tool\_configs}. Timeout minimum \textit{code interpreter} E2B dinaikkan dari 180 detik menjadi 300 detik, dan sub-agen analytics diberi \texttt{timeout\_seconds=1800}. Eksekusi memakai \texttt{Sandbox.run\_code()} menggantikan \texttt{commands.run()} yang sudah tidak digunakan.
    \item \textbf{Penjadwalan Mingguan}: Jadwal analitik dikonfigurasi setiap Jumat pukul 17.00 (opsional).
    \item \textbf{Perbaikan Iteratif}: Selama pengembangan saya menangani error 422 pada pengiriman email, inkonsistensi action items, kendala layout PDF, penanganan penerima, pengambilan detail wajib, dan perbaikan tautan Drive agar hasil sesuai PRD.
\end{itemize}

\subsubsection{DE-PM: Migrasi BOSA ke GLConnectors}

Mulai awal Juni, saya merefaktor tool-tool DE-PM yang masih memakai sesi MCP BOSA ke pola bersama GLConnectors SDK:

\begin{itemize}
    \item \texttt{send\_filtered\_email\_tool} (PR \#866), \texttt{create\_filtered\_drive\_permission\_tool} (PR \#878), dan \texttt{GoogleDocsToHtmlTool} (PR \#892) diubah agar memanggil SDK GLConnectors secara langsung.
    \item Pembungkus \texttt{gl\_connectors\_execute} dihapus, konfigurasi disuntikkan lewat \texttt{get\_tool\_config}, nama variabel \texttt{bosa} diganti \texttt{gl\_connectors}, dan konfigurasi yang sudah tidak dipakai (suffix \texttt{GL\_CONNECTORS\_API\_KEY}/\texttt{TOKEN}) dibersihkan. Refaktor ini diuji hingga cakupan branch 100\%.
    \item Setelah migrasi, data E2B Analytics juga dipindahkan dari Meemo MCP ke GL Connectors.
    \item \textbf{Structured Logging}: Sebagai perbaikan terkait, saya menambahkan modul logging berbasis ContextVar, menghubungkannya ke titik masuk tool runtime, dan memigrasikan tool yang tersisa ke \texttt{ToolConfigMixin} agar \texttt{get\_tool\_config} dan context setup selalu berjalan di dalam try block.
\end{itemize}

\subsection{Refleksi}

Dari membangun pipeline DE-WRI dari nol, saya belajar bahwa sebagian besar bug tidak muncul di unit test, melainkan di interaksi antar komponen. Dua contoh yang paling saya ingat: error runtime AIP yang memaksa TypedDict diganti Pydantic BaseModel, dan \texttt{issue\_id} yang tertinggal dari tuple pengelompokan mingguan sehingga isu berulang tidak terdeteksi. Keduanya hanya muncul saat pipeline benar-benar dijalankan di server, bukan saat diuji lokal.

Pekerjaan migrasi BOSA ke GLConnectors mengajarkan bahwa pola bersama (shared pattern) bernilai ketika jumlah tool banyak; tanpa konsistensi \texttt{get\_tool\_config} dan penamaan variabel, setiap tool akan punya cara sendiri dalam mengelola koneksi dan konfigurasi. Menulis test hingga cakupan 100\% pada tool yang memanggil API eksternal membutuhkan desain mock yang jujur, supaya yang diuji adalah logika tool, bukan perilaku mock itu sendiri.

\clearpage
